\subsection{What the lattice gives without the averaged condition}
\label{ssec:rEunconditional}

The averaged condition $\mathrm{BP}(n_k,\varpi)$ of
Theorem~\ref{thm:avgtransfer} enters Sec.~\ref{ssec:limitE} at one place
only, the decay $e^{-m_*a_kt}$ of Lemma~\ref{lem:limitElattice}(iii), that
is, through the lattice gap of Corollary~\ref{cor:avggap}
(Remark~\ref{rem:limitEscope}). This subsection records what remains when
that gap is removed: the torus limits exist for every coupling, they are
invariant and reflection positive, their two-point functions have
renormalised continuum limits along subsequences with a positive Hamiltonian
on the one-field space, and the short-distance law (AF) alone supplies the
two inputs (SD) and (NE). It then states the missing step as one lemma and
proves that its hypotheses are necessary as well as sufficient.

Throughout, $\beta_k\ge0$ is an arbitrary sequence of couplings; nothing is
assumed about $\mathrm{BP}$, about $D(n,\alpha)$ or about the tuning. For a
state $\omega$ as in Appendix~\ref{app:transfer} with transfer matrix
$\mathbb T$ and vacuum $\Omega$ put
\[
\gamma(\omega)=\bigl\|\mathbb T|_{\Omega^\perp}\bigr\|\in[0,1],
\]
the norm of $\mathbb T$ on the orthogonal complement of the vacuum; since
$\mathbb T$ is self-adjoint and $\mathbb T\Omega=\Omega$, the subspace
$\Omega^\perp$ is invariant under $\mathbb T$. We write $\gamma_k=\gamma(
\omega_k)$. A \emph{lattice gap} $m_{\mathrm{lat}}>0$ means $\gamma_k\le
e^{-m_{\mathrm{lat}}}$ for all $k$; the \emph{physical gap} at level $k$ is
$a_k^{-1}\log(1/\gamma_k)\in[0,\infty]$.

\begin{lemma}[torus limits exist and are reflection positive]\label{rE:torus}
Let $G$ be a compact group with a faithful unitary representation, $\beta\ge0$,
and let $\mu_L$ be the Wilson measure of $\Lambda_L=(\Z/L\Z)^4$ at $\beta$.

\textup{(i)} Every sequence of even integers $L\to\infty$ has a subsequence
along which $\omega(O)=\lim_L\mu_L(O)$ exists for every continuous function
$O$ of finitely many links of $\Z^4$, each $O$ being evaluated on $\Lambda_L$
through a box centred at the origin once $L$ is large. The functional
$\omega$ is a state, given by a probability measure on the configurations of
$\Z^4$.

\textup{(ii)} $\omega$ is invariant under the translations of $\Z^4$ and under
the hypercubic group $B_4$.

\textup{(iii)} $\omega$ is reflection positive for the reflections in
coordinate hyperplanes through sites and in coordinate hyperplanes bisecting
links: $\omega(\vartheta_\mu(\bar F)F)\ge0$ for every $\mu$ and every
continuous gauge-invariant function $F$ of finitely many links with both
endpoints on the positive side of the plane of $\vartheta_\mu$.

\textup{(iv)} The reconstruction of Appendix~\ref{app:transfer} applies to
$\omega$ and gives $\mathcal H_\omega$, $\Omega$, $F\mapsto\Psi_F$ and a
self-adjoint $\mathbb T$ with $0\le\mathbb T\le1$, $\mathbb T\Omega=\Omega$
and $\ip{\Psi_F}{\mathbb T^t\Psi_H}=\omega(\vartheta(\bar F)\tau_tH)$ for
integers $t\ge0$. For $F$ with $\omega(F)=0$, $\Psi_F\perp\Omega$ and
$\|\mathbb T^t\Psi_F\|\le\gamma(\omega)^t\|\Psi_F\|$.
\end{lemma}

\begin{proof}
(i) Let $E_m$ be the set of links of $\Z^4$ with both endpoints in
$[-m,m]^4$. For $L>2m+2$ the quotient map $\Z^4\to\Lambda_L$ is injective on
$[-m,m]^4$, two of its points differing by less than $L$ in every
coordinate, so it identifies $E_m$ with a set of links of $\Lambda_L$ and a
function of the links in $E_m$ with a function on the configurations of
$\Lambda_L$; this defines $\mu_L(O)$ for $O\in C(G^{E_m})$ and $L>2m+2$. The
image measure $\nu_{L,m}$ of $\mu_L$ on $G^{E_m}$ is a Borel probability
measure on a compact metrisable space, so by the Riesz representation
theorem and the Banach--Alaoglu theorem every sequence of them has a weakly
convergent subsequence, and a diagonal argument over $m\in\mathbb N$ gives a
subsequence $(L_j)$ with $\nu_{L_j,m}\to\nu_m$ weakly for every $m$. The
projection $G^{E_{m+1}}\to G^{E_m}$ is continuous and carries $\nu_{L,m+1}$
to $\nu_{L,m}$, hence $\nu_{m+1}$ to $\nu_m$. The Kolmogorov extension
theorem on the countable product $G^{E(\Z^4)}$ gives a probability measure
$\nu$ with marginals $\nu_m$, and $\omega(O)=\int O\,d\nu=\lim_j\mu_{L_j}(O)$
for $O\in C(G^{E_m})$ and every $m$.

(ii) Let $v\in\Z^4$ and $O\in C(G^{E_m})$. The translate $\tau_vO$ lies in
$C(G^{E_{m'}})$ with $m'=m+|v|_\infty$. For $L>2m'+2$ the translation of
$\Lambda_L$ by the class of $v$ is a bijection of its links that preserves
the set of plaquettes and therefore $\mu_L$, and it carries the embedding of
$E_m$ to that of $E_m+v$; hence $\mu_L(\tau_vO)=\mu_L(O)$, and the limit
gives $\omega(\tau_vO)=\omega(O)$. The hypercubic group acts on $\Lambda_L$
fixing the origin, permutes the plaquettes, and preserves $\ReTr Q_p$ under
reversal of orientation because $\ReTr U^{-1}=\ReTr U$ for unitary $U$; it
carries each link variable to the variable of the image link or to its
inverse, and the product Haar measure is invariant under $U\mapsto U^{-1}$
on every link; so $\mu_L$ is $B_4$-invariant, and so is $\omega$.

(iii) Let $F$ be a continuous gauge-invariant function of the links in
$E_m$ lying on the positive side of the plane of $\vartheta_\mu$. For even
$L>2m+2$ the reflection $\vartheta_\mu$ induces a reflection of $\Lambda_L$,
in the plane $\{x_\mu=0\}$ (and its antipode $\{x_\mu=L/2\}$) for the
reflection through sites and in the plane $\{x_\mu=\frac12\}$ for the
reflection bisecting links, the embedded support of $F$ lies strictly on one
side of it, and Osterwalder and Seiler proved $\mu_L(\vartheta_\mu(\bar F)
F)\ge0$ for the Wilson measure of the torus with $\beta\ge0$ for both kinds
of reflection~\cite{OSei78,Sei82}. This is one inequality for one local
function, and it passes to the limit along $(L_j)$.

(iv) By (ii) and (iii), $\omega$ is a state of the kind treated in
Appendix~\ref{app:transfer}, which gives the reconstruction. If
$\omega(F)=0$ then $\ip\Omega{\Psi_F}=\omega(F)=0$, so $\Psi_F\in
\Omega^\perp$, and $\|\mathbb T^t\Psi_F\|\le\gamma(\omega)^t\|\Psi_F\|$
because $\Omega^\perp$ is invariant under $\mathbb T$.
\end{proof}

The proof of Corollary~\ref{cor:avggap} is this lemma followed by the
decay of Theorem~\ref{thm:avgdecay}; the lemma is the part that needs no
hypothesis on $\beta$.

\emph{The gap-free two-point theory.} Let $\omega_k$ be torus limits at
$\beta_k\ge0$ as in Lemma~\ref{rE:torus}, with $\mathbb T_k$, $\Omega_k$ and
$\gamma_k$, and let $\Phi$, $r$, $A_k(f)$, $c_k$, $\mathcal N_k(\varepsilon)$
be as in Sec.~\ref{ssec:limitE}. For $\eta\ge0$ and a direction $\mu$ let
$\mathcal V^\eta_\mu$ be the set of $g\in C_c(\R^4)$ with
$\operatorname{supp}g\subset\{x_\mu>\eta\}$, and for $g\in\mathcal V^0_0$
put
\[
R_k(g)=\frac{\omega_k\bigl(\vartheta(\overline{A_k(g)})A_k(g)\bigr)}
{\mathcal N_k(\varepsilon)}
\]
when $\mathcal N_k(\varepsilon)>0$. Let $\tau_yg=g(\cdot-y)$ for $y\in\R^4$.

\begin{theorem}[the two-point function without the averaged condition]
\label{rE:two}
Let $\varepsilon>0$ and let $k$ be so large that $(r+1)a_k\le\varepsilon$.

\textup{(i)} If $\mathcal N_k(\varepsilon)=0$, then $\omega_k(A_k(f)A_k(g))=0$
for every $(f,g)\in\mathcal C_2(Q,\delta)$ with $\delta\ge8\varepsilon$.
Consequently every sequence of $k$ has a subsequence along which either
$\mathcal N_k(\varepsilon)>0$ for every $k$, or the connected two-point
function $G_k$ of $\Phi$ vanishes, for every $k$, at all $z\in\Z^4$ with
$a_k|z|_\infty>8\varepsilon$.

In \textup{(ii)}--\textup{(vi)} suppose $\mathcal N_k(\varepsilon)>0$ for the
$k$ considered and put $Z_k=\mathcal N_k(\varepsilon)^{-1/2}$.

\textup{(ii)} Let $c\in\R$, $d\ge0$, $f$ supported in a cube of side $L$ in
$\{x_\mu\le c-2\varepsilon\}$ and $g$ in a cube of side $L'$ in
$\{x_\mu\ge c+2\varepsilon+d\}$. Then
\[
Z_k^2\bigl|\omega_k\bigl(A_k(f)A_k(g)\bigr)\bigr|\le\mathsf K_L\mathsf
K_{L'}\,\gamma_k^{\lfloor d/a_k\rfloor}\,\|f\|_\infty\|g\|_\infty ,
\]
and for $(f,g)\in\mathcal C_2(Q,\delta)$ with $Q$ of side $L$ and
$\delta\ge8\varepsilon$,
\[
Z_k^2\bigl|\omega_k\bigl(A_k(f)A_k(g)\bigr)\bigr|\le\mathsf K_{4L/\delta}^2
\mathsf K_{\delta/4}^2\,\|f\|_\infty\|g\|_\infty .
\]

\textup{(iii)} Every sequence of $k$ has a subsequence along which
$S_2(f,g)=\lim_kZ_k^2\omega_k(A_k(f)A_k(g))$ exists for all $f,g\in
C_c(\R^4)$ with $\dist_\infty(\operatorname{supp}f,\operatorname{supp}g)>
8\varepsilon$. It is symmetric and bilinear, bounded as in \textup{(ii)},
invariant under the translations of $\R^4$ and under $B_4$,
$S_2(\bar f,\bar g)=\overline{S_2(f,g)}$, and for every $\mu$
\[
\sum_{i,j}\bar c_ic_j\,S_2(\bar g_i\circ\theta_\mu,g_j)\ge0\qquad
(c_i\in\C,\ g_i\in\mathcal V^{4\varepsilon}_\mu).
\]
It is given by a unique distribution on the open set $\{(x,y)\in\R^8:
|x-y|_\infty>8\varepsilon\}$.

\textup{(iv)} Let $S_2$ be a limit of \textup{(iii)} and $\mu$ a direction.
The form $\ip g{g'}_\mu=S_2(\bar g\circ\theta_\mu,g')$ on
$\mathcal V^{4\varepsilon}_\mu$ is Hermitian and positive semidefinite; let
$\mathcal H^{(1)}_\mu$ be the completion of its quotient by its null space
and $g\mapsto[g]$ the quotient map. There are a self-adjoint operator
$H^{(1)}_\mu\ge0$ on $\mathcal H^{(1)}_\mu$ and a strongly continuous unitary
representation $U_\mu$ of the hyperplane $e_\mu^\perp$ commuting with it
such that, for all $g,g'\in\mathcal V^{4\varepsilon}_\mu$, $t\ge0$ and
$y\perp e_\mu$,
\[
S_2\bigl(\bar g\circ\theta_\mu,\tau_{te_\mu+y}g'\bigr)=\bigl\langle[g],
e^{-tH^{(1)}_\mu}U_\mu(y)[g']\bigr\rangle .
\]
The map $[g]\mapsto[g\circ\pi_\mu]$ is a unitary operator from
$\mathcal H^{(1)}_0$ onto $\mathcal H^{(1)}_\mu$ that carries $H^{(1)}_0$ to
$H^{(1)}_\mu$. If $S_2$ is defined on all separated pairs, the same holds
with $\mathcal V^{4\varepsilon}_\mu$ replaced by $\mathcal V^0_\mu$.

\textup{(v)} For a limit $S_2$ of \textup{(iii)} the following are
equivalent: $S_2$ does not vanish identically; $S_2(\bar g\circ\theta_0,g)>0$
for some $g\in\mathcal V^{4\varepsilon}_0$; $\mathcal H^{(1)}_0\ne0$. Some
limit of \textup{(iii)} does not vanish identically if and only if
$\limsup_kR_k(g)>0$ for some $g\in\mathcal V^{4\varepsilon}_0$, and if
$\liminf_kR_k(g_0)>0$ for one $g_0\in\mathcal V^{4\varepsilon}_0$, then no
limit of \textup{(iii)} vanishes identically. If $S_2$ is defined on all
separated pairs, the equivalence holds with $\mathcal V^{4\varepsilon}_0$
replaced by $\mathcal V^0_0$.

\textup{(vi)} If there is a lattice gap, $\gamma_k\le e^{-m_{\mathrm{lat}}}$
for a fixed $m_{\mathrm{lat}}>0$ and all $k$, then every limit of
\textup{(iii)} vanishes identically.
\end{theorem}

The theorem is Theorem~\ref{thm:limitEtwo}(i)--(iii),(v) and the
constructions of Corollary~\ref{cor:limitEtwo}(b) with the gap removed:
$\gamma_k$ is only the norm of $\mathbb T_k$ on $\Omega_k^\perp$, no bound
on it is assumed, and the Hamiltonian $H^{(1)}_\mu$ is positive but no
lower bound on its spectrum is claimed.

\begin{proof}
We use the following parts of Sec.~\ref{ssec:limitE}, none of which uses
the averaged condition: Lemma~\ref{lem:limitElattice}(i) and (ii), which are
identities; the first two assertions of Lemma~\ref{lem:limitElattice}(iii),
$\omega_k(\vartheta_\mu(\bar F)F)\ge0$ and $0\le\omega_k(\vartheta_\mu(\bar
F)\tau^\mu_{2t}F)\le\omega_k(\vartheta_\mu(\bar F)F)$, whose proof uses only
the reconstruction and $0\le\mathbb T_k\le1$; the identity
$N_\mu(\tau_vF)=N_\mu(F)$ for $v_\mu=0$ from the same proof; and
Lemma~\ref{lem:limitEbox}, whose proof uses Lemma~\ref{lem:limitElattice}(i),
(ii) and $\|\mathbb T_k^{n_0}\|\le1$. All of these hold for the states of
Lemma~\ref{rE:torus}.

(i) Suppose $\mathcal N_k(\varepsilon)=0$. Then $\Psi_{A_k(h)}=0$ for every
$h$ supported in $B_\varepsilon$, and Lemma~\ref{lem:limitEbox} gives
$\Psi_{A_k(g)}=0$ for every $g$ supported in a cube and in
$\{x_0\ge\varepsilon\}$, and, in the form with $\vartheta_\mu$, the same for
$g$ supported in $\{x_\mu\ge\varepsilon\}$ with the vectors of the
reconstruction in the direction $\mu$. Let $f$ be supported in $\{x_\mu\le
c-2\varepsilon\}$ and $g$ in $\{x_\mu\ge c+2\varepsilon\}$, both in cubes. The
first three sentences of the proof of Theorem~\ref{thm:limitEtwo}(i) use only
the invariance of $\omega_k$ under $\pi_\mu$ and the translations and
Lemma~\ref{lem:limitElattice}(i), and give $\omega_k(A_k(f)A_k(g))=
\ip{\Psi_{A_k(h)}}{\mathbb T_k^n\Psi_{A_k(g')}}$ with $h$ and $g'$ supported
in $\{x_0\ge\varepsilon\}$; so $\omega_k(A_k(f)A_k(g))=0$. For $(f,g)\in
\mathcal C_2(Q,\delta)$ with $\delta\ge8\varepsilon$, the decomposition in
the proof of Theorem~\ref{thm:limitEtwo}(ii), with $s=\delta/4$, writes
$\omega_k(A_k(f)A_k(g))$ as a finite sum of such terms, each with pieces in
cubes of side $s$ whose projections on some axis $\mu$ are at distance at
least $\delta/2\ge4\varepsilon$; so the sum vanishes. Given $x,y\in\Z^4$ with
$a_k|x-y|_\infty>8\varepsilon$, choose $f,g\in C_c(\R^4)$ with values in
$[0,1]$, supported in $\ell^\infty$-balls of radius $\rho<a_k/2$ about
$a_kx$ and $a_ky$ with $2\rho\le a_k|x-y|_\infty-8\varepsilon$, and with
$f(a_kx)=g(a_ky)=1$; then $A_k(f)=a_k^4(\Phi_x-c_k)$,
$A_k(g)=a_k^4(\Phi_y-c_k)$, the pair lies in $\mathcal C_2(Q,8\varepsilon)$
for a cube $Q$ containing both balls, and by translation invariance
$a_k^8G_k(y-x)=\omega_k(A_k(f)A_k(g))=0$. The dichotomy follows by passing
to a subsequence on which the sign of $\mathcal N_k(\varepsilon)$ is
constant.

(ii) The proof of Theorem~\ref{thm:limitEtwo}(i) applies verbatim: the
bound $|\ip{\Psi_{A_k(h)}}{\mathbb T_k^n\Psi_{A_k(g')}}|\le\gamma_k^n
\|\Psi_{A_k(h)}\|\|\Psi_{A_k(g')}\|$ holds because both vectors are
orthogonal to $\Omega_k$, the packets being centred, and $\Omega_k^\perp$ is
invariant under $\mathbb T_k$ (Lemma~\ref{rE:torus}(iv)); the norms are
bounded by Lemma~\ref{lem:limitEbox}. The second bound is the proof of
Theorem~\ref{thm:limitEtwo}(ii), which uses the first with $\gamma_k^{\lfloor
d/a_k\rfloor}\le1$.

(iii) Let $\mathcal G$ be the countable set of the proof of
Theorem~\ref{thm:limitEsep}(i). By (ii) the sequence $Z_k^2\omega_k(A_k(g)
A_k(g'))$ is bounded for each pair $g,g'\in\mathcal G$ at distance greater
than $8\varepsilon$, and a diagonal subsequence makes all of them converge.
For $f,g$ at distance $\delta>8\varepsilon$, choose $g_1,g_2\in\mathcal G$
as in that proof with $\|f-g_1\|_\infty,\|g-g_2\|_\infty<\eta$ and supports
within $\eta$ of those of $f,g$, where $\eta<(\delta-8\varepsilon)/4$; the
pairs $(f-g_1,g)$, $(g_1,g-g_2)$ and $(g_1,g_2)$ lie in one $\mathcal C_2(Q',
\delta')$ with $\delta'=\delta-2\eta>8\varepsilon$, and the second bound of
(ii) shows that $Z_k^2\omega_k(A_k(f)A_k(g))$ is a Cauchy sequence. Symmetry
and bilinearity pass to the limit, and so does the bound. Invariance under
dyadic translations and under $B_4$, and the conjugation identity, follow
from Lemma~\ref{lem:limitElattice}(i) and the invariance of $\omega_k$, and
invariance under real translations from the bound and uniform continuity,
exactly as in the proof of Theorem~\ref{thm:limitEsep}(ii), the translated
pairs staying at distance greater than $8\varepsilon$. For the positivity,
let $g=\sum_ic_ig_i$ with $g_i\in\mathcal V^{4\varepsilon}_\mu$; then
$g$ is supported in $\{x_\mu\ge4\varepsilon+\eta'\}$ for some $\eta'>0$. For
$(r+1)a_k\le\varepsilon$ Lemma~\ref{lem:limitElattice}(ii),(iii) gives
$\omega_k(\vartheta_\mu(\overline{A_k(g)})A_k(g))\ge0$, and by
Lemma~\ref{lem:limitElattice}(i) this number equals $\omega_k(A_k(\bar g
\circ\theta_{\mu,k})A_k(g))$. The function $\bar g\circ\theta_{\mu,k}$ is
supported in $\{x_\mu\le-4\varepsilon-\eta'+a_k\}$, so the pairs
$(\bar g\circ\theta_{\mu,k},g)$ and $(\bar g\circ\theta_\mu,g)$ lie in one
$\mathcal C_2(Q',8\varepsilon+\eta')$ for large $k$, and $\|\bar g\circ
\theta_{\mu,k}-\bar g\circ\theta_\mu\|_\infty\to0$ because $\theta_{\mu,k}$
is $\theta_\mu$ followed by a translation of length at most $a_k$ and $g$ is
uniformly continuous. By the bound of (ii) the difference of the two moments
tends to $0$, and $S_2(\bar g\circ\theta_\mu,g)=\lim_kZ_k^2\omega_k(A_k(\bar
g\circ\theta_{\mu,k})A_k(g))\ge0$; expanding $g$ gives the stated
inequality. The distribution is constructed from the bound exactly as in the
proof of Theorem~\ref{thm:limitEsep}(i), with the products of open cubes
replaced by products $U_1\times U_2$ of bounded open sets whose closures are
at $\ell^\infty$-distance greater than $8\varepsilon$; these cover the open
set in question, and the localisation argument given there applies.

(iv) Write $\ip\cdot\cdot$ for $\ip\cdot\cdot_\mu$ and $\tau_t$ for
$\tau_{te_\mu}$. For $g,g'\in\mathcal V^{4\varepsilon}_\mu$ and $t\ge0$ the
pair $(\bar g\circ\theta_\mu,\tau_tg')$ has its supports in $\{x_\mu<-4
\varepsilon\}$ and $\{x_\mu>4\varepsilon+t\}$, at distance greater than
$8\varepsilon$, so $S_2$ is defined on it; and $t\mapsto S_2(\bar g\circ
\theta_\mu,\tau_tg')$ is continuous on $[0,\infty)$ by the bound of (ii)
and the uniform continuity of $g'$, the pairs lying in one $\mathcal C_2(Q',
\delta')$ with $\delta'>8\varepsilon$ for $t$ in a bounded set.

\emph{Hermitian symmetry.} Since $\theta_\mu\in B_4$ and $S_2$ is symmetric,
$B_4$-invariant and satisfies the conjugation identity,
$\ip g{g'}=S_2(\bar g\circ\theta_\mu,g')=S_2(\bar g,g'\circ\theta_\mu)=
S_2(g'\circ\theta_\mu,\bar g)=\overline{S_2(\bar g'\circ\theta_\mu,g)}=
\overline{\ip{g'}g}$. Positivity is (iii). So the quotient by the null
space is a pre-Hilbert space and $\mathcal H^{(1)}_\mu$ is its completion.

\emph{Time translations.} We use $\theta_\mu e_\mu=-e_\mu$, hence
$\overline{\tau_tg}\circ\theta_\mu=\tau_{-t}(\bar g\circ\theta_\mu)$, and the
translation invariance of $S_2$. For $g,g'\in\mathcal V^{4\varepsilon}_\mu$
and $t\ge0$,
\[
\ip g{\tau_tg'}=S_2\bigl(\tau_{-t}(\bar g\circ\theta_\mu),g'\bigr)=
S_2\bigl(\overline{\tau_tg}\circ\theta_\mu,g'\bigr)=\ip{\tau_tg}{g'} .
\tag{a}
\]
Next let $t\ge0$ be dyadic and $k$ so large that $n=t/a_k$ is an integer.
By Lemma~\ref{lem:limitElattice}(ii) the function $F=A_k(g)$ depends on
links in $\{x_\mu\ge1\}$, and Lemma~\ref{lem:limitElattice}(iii) gives
$0\le\omega_k(\vartheta_\mu(\bar F)\tau^\mu_{2n}F)\le\omega_k(\vartheta_\mu(
\bar F)F)$. By Lemma~\ref{lem:limitElattice}(i) the middle term is
$\omega_k(A_k(\bar g\circ\theta_{\mu,k})A_k(\tau_{2t}g))$ and the right term
is the same with $\tau_{2t}g$ replaced by $g$. Multiplying by $Z_k^2$ and
passing to the limit, with the approximation of $\theta_{\mu,k}$ by
$\theta_\mu$ as in (iii), gives
\[
0\le S_2(\bar g\circ\theta_\mu,\tau_{2t}g)\le S_2(\bar g\circ\theta_\mu,g)
\qquad(t\ge0\text{ dyadic}),
\]
and by continuity in $t$ for all $t\ge0$. By (a) the middle term is
$\ip{\tau_tg}{\tau_tg}$. Hence $\|[\tau_tg]\|\le\|[g]\|$, so $T(t)[g]=
[\tau_tg]$ is a well defined linear contraction of the quotient, which
extends to a contraction $T(t)$ of $\mathcal H^{(1)}_\mu$; by (a) it is
symmetric, hence self-adjoint; $\ip{[g]}{T(t)[g]}=\|[\tau_{t/2}g]\|^2\ge0$
on a dense set, so $T(t)\ge0$; and $T(t)T(s)=T(t+s)$, $T(0)=1$. For strong
continuity, let $g\in\mathcal V^{4\varepsilon}_\mu$ and $t\ge0$; the
function $h=\tau_tg-g$ lies in $\mathcal V^{4\varepsilon}_\mu$, and
\[
\|T(t)[g]-[g]\|^2=S_2(\bar h\circ\theta_\mu,h)\le\mathsf K_{4L'/\delta'}^2
\mathsf K_{\delta'/4}^2\|h\|_\infty^2
\]
by the bound of (ii), for a cube of side $L'$ and a separation $\delta'>8
\varepsilon$ that do not depend on $t\in[0,1]$. Since $\|h\|_\infty\to0$ as
$t\to0$, $T(t)\to1$ strongly on a dense set, hence on $\mathcal H^{(1)}_\mu$
because $\|T(t)\|\le1$, and $t\mapsto T(t)$ is strongly continuous on
$[0,\infty)$ by the semigroup law. Now $T(1)$ is injective: $\ker T(2^{-m})=
\ker T(2^{-m-1})$ for every $m$, because $T(2^{-m})=T(2^{-m-1})^2$ with
$T(2^{-m-1})\ge0$; so a vector $v\in\ker T(1)$ lies in $\ker T(2^{-m})$ for
all $m$, and $v=\lim_mT(2^{-m})v=0$. Hence the spectral measure of $T(1)$
has no mass at $0$, and $H^{(1)}_\mu=-\log T(1)$, defined by the functional
calculus on $(0,1]$, is self-adjoint with $H^{(1)}_\mu\ge0$ and $e^{-H^{(1)}
_\mu}=T(1)$. The positive square root of a positive operator is unique, so
$T(2^{-m})=T(1)^{2^{-m}}$ by induction on $m$, and $T(t)=e^{-tH^{(1)}_\mu}$
for every dyadic $t\ge0$; both sides are strongly continuous in $t$, so the
identity holds for all $t\ge0$.

\emph{Spatial translations.} For $y\perp e_\mu$, $\theta_\mu y=y$, so
$\overline{\tau_yg}\circ\theta_\mu=\tau_y(\bar g\circ\theta_\mu)$, and the
translation invariance of $S_2$ gives $\ip{\tau_yg}{\tau_yg'}=\ip g{g'}$;
thus $U_\mu(y)[g]=[\tau_yg]$ defines an isometry of the quotient, with
inverse $U_\mu(-y)$, which extends to a unitary operator; $U_\mu(y)U_\mu(y')
=U_\mu(y+y')$, $U_\mu(y)$ commutes with $T(t)$ because $\tau_y\tau_t=\tau_t
\tau_y$, and $U_\mu$ is strongly continuous by the same estimate as for
$T$. Finally $\ip g{\tau_{te_\mu+y}g'}=\ip{[g]}{T(t)U_\mu(y)[g']}$ by
definition.

\emph{The four directions.} For $g\in\mathcal V^{4\varepsilon}_0$ the
function $g\circ\pi_\mu$ vanishes unless $(\pi_\mu x)_0=x_\mu>4
\varepsilon$, so it lies in $\mathcal V^{4\varepsilon}_\mu$, and the map is
onto because $\pi_\mu$ is an involution. From $\pi_\mu\theta_\mu=\theta_0
\pi_\mu$ and the $B_4$-invariance of $S_2$,
\begin{align*}
S_2\bigl(\overline{g\circ\pi_\mu}\circ\theta_\mu,g'\circ\pi_\mu\bigr)&=
S_2\bigl((\bar g\circ\theta_0)\circ\pi_\mu,g'\circ\pi_\mu\bigr)\\
&=S_2(\bar g\circ\theta_0,g'),
\end{align*}
so $[g]\mapsto[g\circ\pi_\mu]$ is a well defined isometry of the quotients,
which extends to a unitary operator $W_\mu$. Since $\pi_\mu e_0=e_\mu$,
$(\tau_{te_0}g)\circ\pi_\mu=\tau_{te_\mu}(g\circ\pi_\mu)$, so $W_\mu$
intertwines the two semigroups and therefore the two Hamiltonians.

If $S_2$ is defined on all separated pairs, every step above applies with
$4\varepsilon$ replaced by $0$: the pairs used are then separated, the
positivity $\omega_k(\vartheta_\mu(\overline{A_k(g)})A_k(g))\ge0$ holds for
$g$ supported in $\{x_\mu\ge\eta\}$ once $(r+1)a_k\le\eta$, and the
approximation of $\theta_{\mu,k}$ by $\theta_\mu$ keeps the pairs separated.

(v) If $S_2(\bar g\circ\theta_0,g)>0$ for some $g\in\mathcal V^{4
\varepsilon}_0$, then $S_2\ne0$ and $[g]\ne0$. Conversely suppose
$S_2(\bar g\circ\theta_0,g)=0$ for all $g\in\mathcal V^{4\varepsilon}_0$.
By (iv) with $\mu=0$ and $B_4$-invariance, $S_2(\bar g\circ\theta_\mu,g)=
S_2(\overline{g\circ\pi_\mu}\circ\theta_0,g\circ\pi_\mu)=0$ for every
$\mu$ and $g\in\mathcal V^{4\varepsilon}_\mu$, so the positive semidefinite
form $\ip\cdot\cdot_\mu$ vanishes on its diagonal and, by the Cauchy--Schwarz
inequality, identically: $S_2(\bar g\circ\theta_\mu,g')=0$ for $g,g'\in
\mathcal V^{4\varepsilon}_\mu$. With translation invariance this says that
$S_2(f,g)=0$ whenever $f$ is supported in $\{x_\mu<c-4\varepsilon\}$ and $g$
in $\{x_\mu>c+4\varepsilon\}$ for some $c$ and $\mu$, as $f=\bar h\circ
\theta_\mu$ with $h=\bar f\circ\theta_\mu$ after the translation by $-ce_\mu$.
Let $(f,g)\in\mathcal C_2(Q,\delta)$ with $\delta>8\varepsilon$, put
$s=(\delta-8\varepsilon)/4>0$, and write $f=\sum_m\chi^s_mf$, $g=\sum_{m'}
\chi^s_{m'}g$ with the partition of unity of the proof of
Theorem~\ref{thm:limitEsep}(vi), whose members are supported in closed cubes
of side $s$; finitely many terms are nonzero. For a pair of nonzero pieces
pick points $p,q$ of their supports and $\mu$ with $|p_\mu-q_\mu|\ge\delta$;
the projections of the two cubes on the $\mu$ axis are at distance at least
$\delta-2s=\delta/2+4\varepsilon>8\varepsilon$, so the pieces are at
$\ell^\infty$-distance greater than $8\varepsilon$ and lie on opposite sides
of some $c$ with margins greater than $4\varepsilon$. By bilinearity $S_2(f,
g)$ is the sum of the $S_2$ of the pairs of pieces, each of which is $0$.
So $S_2$ vanishes identically, and $\mathcal H^{(1)}_0=0$ means exactly
that $S_2(\bar g\circ\theta_0,g)=0$ for all $g\in\mathcal V^{4\varepsilon}
_0$. This proves the equivalence. If $S_2$ is defined on all separated
pairs, the same argument with $\mathcal V^0_\mu$ in place of
$\mathcal V^{4\varepsilon}_\mu$, pairs at a positive distance $\delta$,
$s=\delta/4$ and margins greater than $0$ gives the equivalence on
$\mathcal V^0_0$.

For $g\in\mathcal V^{4\varepsilon}_0$ and large $k$, Lemma~\ref{lem:limitElattice}(i)
gives $R_k(g)=Z_k^2\omega_k(A_k(\bar g\circ\theta_{0,k})A_k(g))$, and along
any subsequence with a limit $S_2$, the approximation of $\theta_{0,k}$ by
$\theta_0$ in (iii) gives $R_k(g)\to S_2(\bar g\circ\theta_0,g)$. If some
limit $S_2$ is not zero, pick $g$ with $S_2(\bar g\circ\theta_0,g)>0$; then
$\limsup_kR_k(g)>0$. If $\limsup_kR_k(g)>0$, take a subsequence along which
$R_k(g)$ tends to the $\limsup$ and refine it by (iii); the limit satisfies
$S_2(\bar g\circ\theta_0,g)=\limsup_kR_k(g)>0$. If $\liminf_kR_k(g_0)>0$,
every limit has $S_2(\bar g_0\circ\theta_0,g_0)\ge\liminf_kR_k(g_0)>0$.

(vi) Let $(f,g)\in\mathcal C_2(Q,\delta)$ with $\delta>8\varepsilon$ and
decompose as in the proof of the second bound of (ii), with pieces in cubes
of side $s=\delta/4$ whose projections on some axis are at distance at least
$\delta/2$. Each pair of pieces satisfies the hypothesis of the first bound
of (ii) with $d=\delta/2-4\varepsilon>0$, so the lattice quantity is at most
$\mathsf K_{4L/\delta}^2\mathsf K_{\delta/4}^2\|f\|_\infty\|g\|_\infty
e^{-m_{\mathrm{lat}}(d/a_k-1)}$, which tends to $0$.
\end{proof}

The structure that (iii) and (iv) give is, in the language of~\cite{OS73}
restricted to the two-point function at separations greater than
$8\varepsilon$: symmetry, invariance under the translations and the
hypercubic group, and reflection positivity in each of the four coordinate
directions, with a positive Hamiltonian on the one-field space of each
direction. It does not give invariance under $\mathrm O(4)$, the
independence of the limit from the subsequence, a regularity or growth
condition at coinciding points, clustering, or a lower bound on
$\sigma(H^{(1)}_\mu)$; and it does not decide whether $\mathcal H^{(1)}_\mu
\ne0$, which is the content of (v).

\begin{corollary}[the physical gap under (SD)]\label{rE:gapbound}
In the setting of Theorem~\textup{\ref{rE:two}}, with
$\mathcal N_k(\varepsilon)>0$ for large $k$:

\textup{(i)} for $(r+1)a_k\le\varepsilon'<\varepsilon$,
$\mathcal N_k(\varepsilon)\le\mathsf K^2\gamma_k^{2\lfloor(\varepsilon-
\varepsilon')/a_k\rfloor}\mathcal N_k(\varepsilon')$;

\textup{(ii)} if \textup{(SD)} holds, then for every $\varepsilon'\in(0,
\varepsilon)$, with $C_{\varepsilon'}=\limsup_k\mathcal N_k(\varepsilon')/
\mathcal N_k(\varepsilon)<\infty$,
\[
\limsup_k\frac1{a_k}\log\frac1{\gamma_k}\le\frac{\log(\mathsf K^2C_{
\varepsilon'})}{2(\varepsilon-\varepsilon')} ;
\]
in particular a lattice gap is incompatible with \textup{(SD)}.

\textup{(iii)} \textup{(SD)} holds if and only if \textup{(UV)} holds, and if
and only if \textup{(M$_{\mathrm s}$)} holds for $n=2$ with the constants
$Z_k=\mathcal N_k(\varepsilon)^{-1/2}$.
\end{corollary}

\begin{proof}
(i) is the first inequality of Proposition~\ref{prop:limitEsd}(i); its
proof uses Lemma~\ref{lem:limitElattice}(i), the identity $\Psi_{A_k(h)}=
\mathbb T_k^n\Psi_{A_k(h')}$ from the proof of
Lemma~\ref{lem:limitElattice}(iii), $\Psi_{A_k(h')}\perp\Omega_k$ for the
centred packet, $\|\mathbb T_k^n\Psi\|\le\gamma_k^n\|\Psi\|$ on
$\Omega_k^\perp$ (Lemma~\ref{rE:torus}(iv)) and Lemma~\ref{lem:limitEbox};
the second inequality there, $\gamma_k\le e^{-m_*a_k}$, is the only use of
the averaged condition and is not needed. (ii) is the proof of
Proposition~\ref{prop:limitEsd}(iii) from (i): with $n=\lfloor(\varepsilon-
\varepsilon')/a_k\rfloor\ge(\varepsilon-\varepsilon')/a_k-1$ and $C>C_{
\varepsilon'}$, (i) gives $\gamma_k^{-2n}\le\mathsf K^2C$ for large $k$,
that is $a_k^{-1}\log(1/\gamma_k)\le\log(\mathsf K^2C)/(2(\varepsilon-
\varepsilon'-a_k))$, and $C\downarrow C_{\varepsilon'}$. Under a lattice gap
the left side is at least $m_{\mathrm{lat}}/a_k\to\infty$. (iii) The
equivalence with (UV) is Proposition~\ref{prop:limitEsd}(ii), whose proof
uses only the commutation of the spectral projections with $\mathbb T_k$,
$\|\mathbb T_k\|\le1$, Lemma~\ref{lem:limitEbox} and spectral measures; the
equivalence with (M$_{\mathrm s,2}$) is Corollary~\ref{cor:limitEtwo}(a),
whose proof uses Theorem~\ref{thm:limitEtwo}(i),(ii) with $\varepsilon'$ in
place of $\varepsilon$, that is, Theorem~\ref{rE:two}(ii).
\end{proof}

\emph{The short-distance law without the averaged condition.} The
hypotheses of Theorem~\ref{thm:limitEaf} are the averaged condition and
(AF). The averaged condition can be dropped from parts (i)--(iv), with one
change in the proof of (iv): (NE) is read off the lower bound of (AF)
directly.

\begin{theorem}[(AF) alone gives the two-point inputs]\label{rE:af}
Let $\omega_k$ be torus limits at arbitrary $\beta_k\ge0$ and assume
\textup{(AF)} for constants $Z_k$. Then:

\textup{(i)} the bound of Theorem~\textup{\ref{thm:limitEaf}(i)} holds, so
\textup{(M$_{\mathrm s,2}$)} holds for the constants $Z_k$, every limit $S_2$
of $Z_k^2\omega_k(A_k(f)A_k(g))$ along a subsequence satisfies
$|S_2(f,g)|\le C_1\iint|f(x)||g(y)|w_\gamma(|x-y|)\,dx\,dy$ on the separated
pairs, and extends uniquely to a continuous linear functional on
$\mathcal S_0(\R^8)$;

\textup{(ii)} for $f,g\ge0$ forming a separated pair with $|x-y|\le r_0$ for
$x\in\operatorname{supp}f$, $y\in\operatorname{supp}g$, every such limit has
$S_2(f,g)\ge c_1\iint f(x)g(y)w_\gamma(|x-y|)\,dx\,dy$; in particular no
limit vanishes identically;

\textup{(iii)} for every $\varepsilon>0$, $\mathcal N_k(\varepsilon)>0$ for
large $k$ and $0<\liminf_kZ_k^2\mathcal N_k(\varepsilon)\le\limsup_kZ_k^2
\mathcal N_k(\varepsilon)<\infty$;

\textup{(iv)} \textup{(SD$_{\mathrm p}$)} holds with $p_0=4$, hence
\textup{(SD)} and \textup{(UV)}; \textup{(NE)} holds; the physical gap
$a_k^{-1}\log(1/\gamma_k)$ is bounded in $k$; and for the constants
$Z^\sharp_k=\mathcal N_k(\varepsilon)^{-1/2}$ every limit $S^\sharp_2$ of
$(Z^\sharp_k)^2\omega_k(A_k(f)A_k(g))$ along a subsequence exists on all
separated pairs, is a positive multiple of a limit of \textup{(i)}, does not
vanish identically, extends to $\mathcal S_0(\R^8)$, and has
$\mathcal H^{(1)}_\mu\ne0$ for every $\mu$;

\textup{(v)} \textup{(AF)} fails for every choice of $Z_k$ if there is a
lattice gap, in particular under the hypotheses of
Corollary~\textup{\ref{cor:limitEstrong}} or
\textup{\cite[Corollary~322]{BBLR26}}.
\end{theorem}

\begin{proof}
(i) and (ii) are proved exactly as Theorem~\ref{thm:limitEaf}(i),(ii): the
proofs use Lemma~\ref{lem:limitElattice}(i), translation invariance, the two
bounds of (AF), the counting of lattice points, the convergence of Riemann
sums, and Lemma~\ref{lem:limitEextend}; the averaged condition is not used.
In particular (M$_{\mathrm s,2}$) holds, so the diagonal argument of the
proof of Theorem~\ref{thm:limitEsep}(i) gives limits on all separated pairs
along subsequences, and the proofs of Theorem~\ref{thm:limitEsep}(ii) and
(iii) for $n=2$, which use only (M$_{\mathrm s,2}$),
Lemma~\ref{lem:limitElattice}(i),(ii) and the positivity part of
Lemma~\ref{lem:limitElattice}(iii), show that every such limit is
symmetric, invariant under the translations and $B_4$, satisfies the
conjugation identity, and has $S_2(\bar g\circ\theta_\mu,g)\ge0$ for
$g\in\mathcal V^0_\mu$. Theorem~\ref{rE:two}(iv), in its form for limits
defined on all separated pairs, applies to these limits with the constants
$Z_k$: its proof used of the constants only that $S_2$ is a limit of
$Z_k^2\omega_k(A_k(\cdot)A_k(\cdot))$ with the stated invariance,
positivity and boundedness, the bound of Theorem~\ref{rE:two}(ii) being
replaced by (M$_{\mathrm s,2}$).

(iii) The upper bound is as in the proof of Theorem~\ref{thm:limitEaf}(iii):
with $Q_0$ the cube of side $L_0=2(\varepsilon+2+2a_0)$ there, for $h$
supported in $B_{\varepsilon'}$ with $\|h\|_\infty\le1$, $\varepsilon'\le
\varepsilon$ and $a_k\le\varepsilon'$, the pair $(\bar h\circ\theta_{0,k},h)$
lies in $\mathcal C_2(Q_0,\min(\varepsilon',1))$, so (i) gives
\[
Z_k^2\mathcal N_k(\varepsilon')\le C_2\max(1,\varepsilon'^{-4}),\qquad
C_2=100\,C_1(1+L_0)^4 ,
\tag{b}
\]
for large $k$. For the lower bound suppose $Z_k^2\mathcal N_k(\varepsilon)
\to0$ along a subsequence, and refine it to a limit $S_2$ on all separated
pairs. Let $g\in\mathcal V^0_0$ be supported in a cube of side $L$ and in
$\{x_0\ge\eta\}$, $\eta>0$, and let $t\ge4\varepsilon+1$ be dyadic. For
large $k$ the function $\bar g\circ\theta_{0,k}$ is supported in $\{x_0\le
-\eta+a_k\}$ and $\tau_{te_0}g$ in $\{x_0\ge\eta+t\}$; with $c=2\varepsilon
-\eta+a_k$ and $d=0$ these are the supports of the first bound of
Theorem~\ref{rE:two}(ii), whose proof does not use $\mathcal N_k(
\varepsilon)>0$ once it is multiplied by $\mathcal N_k(\varepsilon)$, and
that bound gives
\[
Z_k^2\bigl|\omega_k\bigl(A_k(\bar g\circ\theta_{0,k})A_k(\tau_{te_0}g)\bigr)
\bigr|\le\mathsf K_L^2\,Z_k^2\mathcal N_k(\varepsilon)\,\|g\|_\infty^2\to0 .
\]
By the approximation of $\theta_{0,k}$ by $\theta_0$, $S_2(\bar g\circ
\theta_0,\tau_{te_0}g)=0$. By Theorem~\ref{rE:two}(iv) for this $S_2$ on
$\mathcal V^0_0$, this number is $\|e^{-tH^{(1)}_0/2}[g]\|^2$, and
$e^{-tH^{(1)}_0/2}$ is injective, so $[g]=0$, that is $S_2(\bar g\circ
\theta_0,g)=0$ for every $g\in\mathcal V^0_0$. By the argument of
Theorem~\ref{rE:two}(v), applied with $\mathcal V^0_\mu$ and separated pairs,
$S_2$ vanishes identically. This contradicts (ii). Hence $\liminf_kZ_k^2
\mathcal N_k(\varepsilon)>0$, and in particular $\mathcal N_k(\varepsilon)>0$
for large $k$.

(iv) Let $c_\varepsilon=\frac12\liminf_kZ_k^2\mathcal N_k(\varepsilon)>0$.
By (b) and (iii), for $\varepsilon'\in(0,\varepsilon]$ and large $k$,
\[
\mathcal N_k(\varepsilon')\le\frac{C_2\max(1,\varepsilon'^{-4})}{c_\varepsilon}
\mathcal N_k(\varepsilon)\le C_0\Bigl(\frac\varepsilon{\varepsilon'}\Bigr)^4
\mathcal N_k(\varepsilon)
\]
with $C_0=C_2\max(1,\varepsilon^{-4})/c_\varepsilon$, which is (SD$_{\mathrm p}$) with $p_0=4$; it gives (SD), and (UV) and the
bound on the physical gap by Corollary~\ref{rE:gapbound}.

For (NE) let $g_0\in C_c(\R^4)$ take values in $[0,1]$, be supported in
$[\frac{r_0}8,\frac{r_0}4]\times[0,\frac{r_0}8]^3$ and be positive on the
interior of this box. Let $x,y\in\Z^4$ with $g_0(a_kx)g_0(a_ky)>0$. The
plane of $\vartheta$ is $\{x_0=c\}$ with $|c|\le\frac12$, so $\vartheta x=
(2c-x_0,x_1,x_2,x_3)$ and the time component of $y-\vartheta x$ is
$x_0+y_0-2c\ge r_0/(4a_k)-1\ge2r+1$ for large $k$, while $|y-\vartheta x|\le
|a_k^{-1}(a_ky-\theta_0(a_kx))|+1\le a_k^{-1}r_0\sqrt{\tfrac14+\tfrac3{64}}
+1\le r_0/a_k$ for large $k$. So both bounds of (AF) apply to every such
pair, and with $p_k(x,y)=a_k^8g_0(a_kx)g_0(a_ky)$ and
$Y_k=Z_k^2\,\omega_k(\vartheta(\overline{A_k(g_0)})A_k(g_0))$,
\begin{align*}
Y_k&=\sum_{x,y}p_k(x,y)\,Z_k^2G_k(y-\vartheta x)\\
&\ge c_1\sum_{x,y}p_k(x,y)\,w_\gamma(a_k|y-\vartheta x|),
\end{align*}
the identity being Lemma~\ref{lem:limitElattice}(i) and translation
invariance. On the support of $g_0\otimes g_0$ the distance $|y-\theta_0x|$
lies in $[\frac{r_0}4,r_0]$, where $w_\gamma$ is Lipschitz, and
$|a_k(y-\vartheta x)|$ differs from $|a_ky-\theta_0(a_kx)|$ by at most
$a_k$; so the right side is a Riemann sum, up to an error $O(a_k)$, of the
continuous compactly supported function $(x,y)\mapsto g_0(x)g_0(y)w_\gamma(
|y-\theta_0x|)$, and
\[
\liminf_kY_k\ge c_1\iint g_0(x)g_0(y)\,w_\gamma(|y-\theta_0x|)\,dx\,dy=:c_1I>0 .
\]
With (b) at $\varepsilon'=\varepsilon$ this gives
\[
\liminf_kR_k(g_0)\ge\frac{c_1I}{C_2\max(1,\varepsilon^{-4})}>0 ,
\]
which is (NE).

By (SD) and Corollary~\ref{rE:gapbound}(iii), (M$_{\mathrm s,2}$) holds for
$Z^\sharp_k$, so the limits $S^\sharp_2$ exist on all separated pairs along
subsequences, and by (SD$_{\mathrm p}$) and
Proposition~\ref{prop:limitEtwotempered}(a),(b), whose extension step uses
only the proof of Corollary~\ref{cor:limitEtwo}(a), that is
Theorem~\ref{rE:two}(ii), and Lemma~\ref{lem:limitEextend}, they extend to
$\mathcal S_0(\R^8)$; the invariance and positivity of the extensions, which
that proposition takes from Theorem~\ref{thm:limitEtempered}(iii) and
Lemma~\ref{lem:limitEcut}(iii), are not needed here. Along a subsequence with a limit $S^\sharp_2$,
refine so that $Z_k^2\mathcal N_k(\varepsilon)\to\lambda\in(0,\infty)$, by
(iii), and so that the limit $S_2$ with the constants $Z_k$ exists; then
$S^\sharp_2=S_2/\lambda$, and $S^\sharp_2$ is not identically zero by (ii).
By Theorem~\ref{rE:two}(v) in its form on $\mathcal V^0_\mu$,
$\mathcal H^{(1)}_0\ne0$, hence $\mathcal H^{(1)}_\mu\ne0$ for every $\mu$
by the unitary equivalence of Theorem~\ref{rE:two}(iv).

(v) Under (AF), (iii) and (iv) give $\mathcal N_k(\varepsilon)>0$ for large
$k$ and (SD), and Corollary~\ref{rE:gapbound}(ii) excludes a lattice gap.
Corollary~\ref{cor:limitEstrong} provides the lattice gap $3.28/(625e^2)$
for $\SU(2)$ and $\beta_{W,k}\le\frac18$, and \cite[Corollary~322]{BBLR26}
the corresponding gap for $N\beta\le\kappa_G/48$.
\end{proof}

The theorem separates the two halves of the problem for $n=2$: the
short-distance law (AF) is a statement about the Wilson measure at lattice
distances between $2r+1$ and $r_0/a_k$ and carries all of (SD), (UV) and
(NE); the averaged condition is a statement about distances of order $\ell$
in physical units and carries the gap $m_*$. Neither implies the other, and
(AF) forces the physical gap to stay bounded while the lattice gap tends to
zero.

\emph{Flowed packets in a periodic spatial box.} For the spatially flowed
observables of Sec.~\ref{ssec:flowlimit} the analogous unconditional
statement is already in the paper; we record it with the inventory of what
it gives, since it is the only place where positivity of the renormalising
norm is a theorem rather than a hypothesis.

\begin{proposition}[flowed packets, unconditional part]\label{rE:flow}
Let $G=\SU(N)$, replace $\Z^4$ by $\Z\times(\Z/L_k\Z)^3$ with $L_k\ge3$, let
$\omega_k$ be the infinite-time Wilson state of the box at an arbitrary
$\beta_k>0$, and let $\Phi$, $\tau$, $f_0=g\otimes h$ with $g,h\ge0$ and
$g(a_kn)h(a_kx)>0$ for some $n\ge1$ and some $x$, $A_k$, $C_k(\sigma)$ be as
in Sec.~\textup{\ref{ssec:flowlimit}} and
Proposition~\textup{\ref{prop:flowpositive}}. Then $C_k(0)>0$ for every $k$,
and with $Z_k=C_k(0)^{-1/2}$ statements \textup{(i)}--\textup{(iv)} of
Theorem~\textup{\ref{thm:flowlimit}} hold with $\mathcal D$ replaced by its
points $(t,0)$: along a subsequence the renormalised flowed packets have a
continuum limit in the sense of Definition~\textup{\ref{def:limit}}, a
Hilbert space with vacuum $\Omega$, a self-adjoint $H\ge0$ with $H\Omega=0$
and vectors $\psi(f)$ with dense range reproducing the limits of the
reflected correlations at dyadic times, together with a spectral measure
$\mu$ on $[0,\infty]$ with $\lim_kC_k(\sigma)/C_k(0)=\int e^{-\sigma E}d\mu(E)$;
the limit is nontrivial if and only if $\mu([0,\infty))>0$, if and only if
$\lim_kC_k(\sigma)/C_k(0)>0$ along the subsequence for one, equivalently
for every, dyadic $\sigma>0$; and some subsequence has a nontrivial limit if
and only if $\limsup_kC_k(\sigma)/C_k(0)>0$ for one dyadic $\sigma>0$. On the infinite
lattice, with $\omega_k$ the torus limits of Lemma~\textup{\ref{rE:torus}}
and $C_k(0)>0$ assumed, the same holds with all of $\mathcal D$, and the
dyadic spatial translations act by unitaries commuting with $H$.
\end{proposition}

\begin{proof}
Proposition~\ref{prop:flowpositive} gives $C_k(0)>0$ in the box, using
L\"uscher's strict positivity of the transfer matrix~\cite{Lus77}. The proof
of Theorem~\ref{thm:flowlimit}(i)--(iv) uses only the reflection positivity
of $\omega_k$, its transfer matrix with $0\le\mathbb T_k\le1$, the spatial
translations where $\mathcal D$ contains them, Proposition~\ref{prop:spatialflow}(ii)
and Appendix~\ref{app:spectral}, as stated after
Proposition~\ref{prop:flowshift}; part (v) is the only part that uses a
gap. On the infinite lattice the states of Lemma~\ref{rE:torus} have these
properties.
\end{proof}

What Proposition~\ref{rE:flow} gives is the Hilbert space, the positive
Hamiltonian, the time semigroup and the reflection-positive inner product
of one renormalised flowed packet and its translates. It does not give
Euclidean invariance, independence of the subsequence, correlations of more
than two renormalised fields, clustering, a gap, or nontriviality; the last
is the ratio condition, which Theorem~\ref{thm:strongratio} and
Proposition~\ref{prop:limitEnorm}(iii) show to fail wherever the lattice
gap stays bounded below, so that it is a statement about the tuning
$\beta_k\to\infty$ alone.

\subsection{The missing step}\label{ssec:rEmissing}

We state the inputs of Sec.~\ref{ssec:limitE} in one place, prove that
they suffice for (E) at separated points, and prove that the two-point
inputs are necessary. Let $\omega_k$ be torus limits along a tuning
$(a_k,\beta_k)$ as in Lemma~\ref{rE:torus} and $\Phi$ as in
Sec.~\ref{ssec:limitE}; the constants $Z_k>0$ are those of each input.

\begin{enumerate}[label=\textup{(I\arabic*)},leftmargin=*]
\item\label{rE:I1} \emph{The averaged condition along the tuning.} There
are $\varpi>0$, $\ell>0$ and integers $n_k\ge2$ with $n_ka_k\le\ell$ such
that $\mathrm{BP}(n_k,\varpi)$ holds for the Wilson measure of $\Lambda_L$
at $\beta_k$ for all $L$ in the infinite set of even integers along which
$\omega_k$ is formed (the hypothesis of Theorem~\ref{thm:avgtransfer}).
\item\label{rE:I2} \emph{The two-sided short-distance law.} (AF) holds for
the constants $Z_k$: for $\gamma\ge0$, $0<c_1\le C_1$, $r_0\in(0,1]$ and all
large $k$, $Z_k^2|G_k(z)|\le C_1w_\gamma(a_k|z|)$ for $|z|_\infty\ge2r+1$, and
$Z_k^2G_k(z)\ge c_1w_\gamma(a_k|z|)$ if moreover $a_k|z|\le r_0$.
\item\label{rE:I3} \emph{Reflected norms of products with a rate.}
(G$_{\mathrm s}$) holds for the constants $Z_k$; by
Proposition~\ref{prop:limitEreflect} this is equivalent, up to the rate, to
(RN$_m$) for every $m\ge1$.
\item\label{rE:I4} \emph{One rotation.} (W$_{\mathrm s}$): a rotation of
infinite order is an asymptotic symmetry of the moments of the separated
tuples.
\item\label{rE:I5} \emph{Uniqueness.} (U$_{\mathrm s}$): all subsequential
limits of the moments of separated tuples coincide.
\end{enumerate}

\begin{theorem}[the inputs suffice]\label{rE:suffices}
\textup{(a)} Under \textup{\ref{rE:I1}} and \textup{\ref{rE:I2}}, with
$Z^\sharp_k=\mathcal N_k(\varepsilon)^{-1/2}$: \textup{(SD$_{\mathrm p}$)},
\textup{(SD)}, \textup{(UV)} and \textup{(NE)} hold; $Z_k/Z^\sharp_k$ is
bounded above and below by positive constants; every limit $S_2$ of
Corollary~\textup{\ref{cor:limitEtwo}(b)} is not identically zero, extends
to $\mathcal S_0(\R^8)$, is invariant under the translations and $B_4$,
reflection positive in every coordinate direction, and its one-field
Hamiltonians $H^{(1)}_\mu$ have spectrum in $[m_*,\infty)$ and not empty,
$m_*=16\varpi/(625e^2\ell)$.

\textup{(b)} Under \textup{\ref{rE:I1}}--\textup{\ref{rE:I5}} with one
sequence of constants $Z_k$, the moments $\omega_k(\prod_jZ_kA_k(f_j))$
converge for every separated tuple; their limits, extended to
$\mathcal S_0(\R^{4n})$ by Theorem~\textup{\ref{thm:limitEtempered}} and with
$S_0=1$, satisfy the linear growth condition of~\cite{OS75} and the axioms
of~\cite{OS73} of Euclidean invariance, reflection positivity, symmetry and
clustering; the Hilbert space and Hamiltonian of their reconstruction are
$\mathcal H_0$ and $H_0$ of Theorem~\textup{\ref{thm:limitEsep}}, with
$\sigma(H_0)\subset\{0\}\cup[m_*,\infty)$, $\ker H_0=\C\Omega_0$ and
$\sigma(H_0)\cap[m_*,\infty)\ne\emptyset$; and they are not the Schwinger
functions of the vacuum alone. That is, \textup{(E)} holds for the field
$\Phi$ at separated points, with a mass gap of at least $m_*$.
\end{theorem}

\begin{proof}
(a) Theorem~\ref{rE:af}(iii),(iv) gives (SD$_{\mathrm p}$), (SD), (UV), (NE),
the comparability of the constants, the extension to $\mathcal S_0(\R^8)$
and the nonvanishing of every limit of Corollary~\ref{cor:limitEtwo}(b);
\ref{rE:I1} is not used for these. Under \ref{rE:I1},
Corollary~\ref{cor:avggap} gives $\gamma_k\le e^{-m_*a_k}$, and
Corollary~\ref{cor:limitEtwo}(b) gives $\sigma(H^{(1)})\subset[m_*,\infty)$;
the invariance, positivity and the unitary equivalence of the directions
are Theorem~\ref{rE:two}(iii),(iv) on $\mathcal V^0_\mu$, and
$\mathcal H^{(1)}_\mu\ne0$ is Theorem~\ref{rE:af}(iv), so the spectrum is
not empty.

(b) \ref{rE:I3} gives (M$_{\mathrm s}$). Corollary~\ref{cor:limitEOS}, whose
hypotheses are \ref{rE:I1}, (G$_{\mathrm s}$), (W$_{\mathrm s}$) and
(U$_{\mathrm s}$), gives convergence, the growth condition, the four axioms,
the identification of the reconstruction with $(\mathcal H_0,H_0)$,
$\sigma(H_0)\subset\{0\}\cup[m_*,\infty)$ and $\ker H_0=\C\Omega_0$. For
nontriviality, Theorem~\ref{rE:af}(ii) shows that $S_2$ is not identically
zero on the separated pairs, so $\mathcal H_0\ne\C\Omega_0$ by
Theorem~\ref{thm:limitEsep}(vi), and $\sigma(H_0)\cap[m_*,\infty)\ne
\emptyset$ by Theorem~\ref{thm:limitEsep}(iv),(vi). The reconstruction
theorem of~\cite{OS75} turns Schwinger functions with these properties into
a Wightman theory on $\mathcal H_0$ with Hamiltonian $H_0$.
\end{proof}

\begin{proposition}[the two-point inputs are necessary]\label{rE:necessary}
Let $\omega_k$ be torus limits along any tuning, $Z_k>0$ arbitrary
constants, and suppose that for every $n$ and every separated tuple the
limit $S_n(f_1,\dots,f_n)=\lim_k\omega_k(\prod_jZ_kA_k(f_j))$ exists along
the full sequence and that $S_2$ is not identically zero. Then:

\textup{(i)} \textup{(M$_{\mathrm s}$)} holds for the constants $Z_k$;

\textup{(ii)} for every $\varepsilon>0$, $\mathcal N_k(\varepsilon)>0$ for
large $k$ and $0<\liminf_kZ_k^2\mathcal N_k(\varepsilon)\le\limsup_kZ_k^2
\mathcal N_k(\varepsilon)<\infty$; so $Z_k$ is comparable to
$\mathcal N_k(\varepsilon)^{-1/2}$;

\textup{(iii)} \textup{(SD)}, \textup{(UV)} and \textup{(NE)} hold for the
constants $\mathcal N_k(\varepsilon)^{-1/2}$;

\textup{(iv)} $\limsup_ka_k^{-1}\log(1/\gamma_k)<\infty$: the physical gap of
the torus limits stays bounded, and in particular the lattice gap tends to
zero. Consequently no tuning along which the torus limits have a lattice
gap, in particular none that stays in the range of
Corollary~\textup{\ref{cor:limitEstrong}} or
\textup{\cite[Corollary~322]{BBLR26}}, admits such limits, for any
choice of $Z_k$.
\end{proposition}

\begin{proof}
(i) Convergent sequences are bounded, so the hypothesis of the converse in
Theorem~\ref{thm:limitEsep} holds. We repeat its proof. Fix $n$, a cube $Q$
and $\delta>0$, and cover $Q$ by the closed cubes of side $s=\delta/4$ that
carry the functions $\chi^s_m$ of the proof of Theorem~\ref{thm:limitEsep}(vi);
at most $(2L/s+3)^4$ of them meet $Q$. For $(f_j)\in\mathcal C_n(Q,\delta)$
every nonzero term $\omega_k(\prod_jZ_kA_k(\chi^s_{m_j}f_j))$ has its $n$
cubes at pairwise distance at least $\delta-2s=\delta/2$. Fix such cubes
$C_1,\dots,C_n$, let $X_j$ be the Banach space of continuous functions on
$\R^4$ supported in $C_j$ with the supremum norm, and let $B_k(h_1,\dots,
h_n)=\omega_k(\prod_jZ_kA_k(h_j))$ on $X_1\times\cdots\times X_n$. Each $B_k$
is a bounded $n$-linear form, since $|A_k(h)|\le2\|\Phi\|_\infty\|h\|_\infty
\sum_xa_k^4\mathbf 1_{C_j}(a_kx)$, and for each point the sequence $(B_k)$
is bounded because it converges. A pointwise bounded family of bounded
$n$-linear forms on a product of Banach spaces is uniformly bounded: for
$n=1$ this is the uniform boundedness principle; for $n\ge2$ and fixed $h_1$
the forms $B_k(h_1,\cdot)$ are bounded $(n-1)$-linear and pointwise bounded,
so by induction $\sup_k\|B_k(h_1,\cdot)\|<\infty$, and the linear maps $h_1
\mapsto B_k(h_1,\cdot)$ into the Banach space of bounded $(n-1)$-linear
forms are pointwise bounded, hence uniformly bounded. Summing the resulting
bounds over the finitely many choices of cubes gives $M_n(Q,\delta)$.

(ii) Upper bound: for $h$ supported in $B_\varepsilon$ with $\|h\|_\infty\le
1$ and $a_k\le\varepsilon$, the pair $(\bar h\circ\theta_{0,k},h)$ lies in
$\mathcal C_2(Q_0,\varepsilon)$ for the cube $Q_0$ of the proof of
Theorem~\ref{rE:af}(iii), and $\omega_k(\vartheta(\overline{A_k(h)})A_k(h))=
\omega_k(A_k(\bar h\circ\theta_{0,k})A_k(h))$; so (i) gives $Z_k^2\mathcal
N_k(\varepsilon)\le M_2(Q_0,\varepsilon)$. Lower bound: by (i) and the proofs
of Theorem~\ref{thm:limitEsep}(ii),(iii) for $n=2$, $S_2$ is symmetric,
invariant, satisfies the conjugation identity and is reflection positive on
$\mathcal V^0_\mu$, so Theorem~\ref{rE:two}(iv) on $\mathcal V^0_\mu$
applies, and by the argument of Theorem~\ref{rE:two}(v) with separated
pairs, $S_2\ne0$ gives $g\in\mathcal V^0_0$ with $[g]\ne0$. Since
$e^{-\varepsilon H^{(1)}_0}$ is injective, $S_2(\overline{\tau_{\varepsilon
e_0}g}\circ\theta_0,\tau_{\varepsilon e_0}g)=\|e^{-\varepsilon H^{(1)}_0}
[g]\|^2>0$, and $g'=\tau_{\varepsilon e_0}g$ is supported in a cube of side
$L$ and in $\{x_0\ge\varepsilon\}$. By Lemma~\ref{lem:limitElattice}(i) and
Lemma~\ref{lem:limitEbox},
\begin{align*}
Z_k^2\,\omega_k\bigl(A_k(\bar g'\circ\theta_{0,k})A_k(g')\bigr)&=Z_k^2\,
\omega_k\bigl(\vartheta(\overline{A_k(g')})A_k(g')\bigr)\\
&\le\mathsf K_L^2\,Z_k^2\mathcal N_k(\varepsilon)\,\|g\|_\infty^2 ,
\end{align*}
and the left side tends to $S_2(\bar g'\circ\theta_0,g')>0$ by the
approximation of $\theta_{0,k}$ by $\theta_0$. Hence $\liminf_kZ_k^2
\mathcal N_k(\varepsilon)>0$.

(iii) For $\varepsilon'\in(0,\varepsilon)$, (ii) at $\varepsilon'$ and at
$\varepsilon$ gives $\mathcal N_k(\varepsilon')\le M_2(Q_0,\varepsilon')
Z_k^{-2}\le C\,\mathcal N_k(\varepsilon)$ for large $k$, which is (SD);
(UV) follows by Corollary~\ref{rE:gapbound}(iii). For (NE), the function
$g'$ of (ii) has $\liminf_kR_k(g')\ge S_2(\bar g'\circ\theta_0,g')/
\limsup_kZ_k^2\mathcal N_k(\varepsilon)>0$.

(iv) is Corollary~\ref{rE:gapbound}(ii) with (SD) from (iii); a lattice gap
$m_{\mathrm{lat}}$ gives $a_k^{-1}\log(1/\gamma_k)\ge m_{\mathrm{lat}}/a_k\to
\infty$. Corollary~\ref{cor:limitEstrong} and~\cite[Corollary~322]{BBLR26}
provide the lattice gaps in their ranges.
\end{proof}

\begin{remark}[independence of the inputs]\label{rE:independent}
The examples of Sec.~\ref{ssec:limitE}, states on $\Z^4$ invariant under
translations and $B_4$ and reflection positive in the four coordinate
directions, show that the consequences of the inputs do not imply one
another. Proposition~\ref{prop:limitEheavy} gives such states with a
physical gap bounded above and below for which (SD) fails, so the lattice
gap $\gamma_k\le e^{-m_*a_k}$ that \ref{rE:I1} provides does not give
\ref{rE:I2}. Theorem~\ref{rE:af}(iv) shows that \ref{rE:I2} bounds the
physical gap above and gives no lower bound on it, and
Proposition~\ref{prop:limitEafexample} gives states satisfying (AF) with
physical gap $\frac12$; whether (AF) can hold with physical gap tending to
zero is not decided by the examples. Proposition~\ref{prop:limitEthree}
gives states with a lattice gap, (SD$_{\mathrm p}$) and (NE) for which
(RN$_m$), $m\ge2$, fails, so a gap and the two-point conditions do not give
\ref{rE:I3}; these states do not satisfy (AF), and no example in
Sec.~\ref{ssec:limitE} decides whether \ref{rE:I1} and \ref{rE:I2} together
give \ref{rE:I3}. Proposition~\ref{prop:limitEfactorial} gives states for
which (M$_{\mathrm s}$) and (RN$_m$) hold for every $m$ while the rate of
(G$_{\mathrm s}$) fails, and Proposition~\ref{prop:limitEnotempered} shows
that the separated-point structure without a rate does not give
temperedness, which is why \ref{rE:I3} carries the rate.
Proposition~\ref{prop:limitEafexample} gives states satisfying (AF),
(G$_{\mathrm s}$) and (U$_{\mathrm s}$) that are not invariant under any
rotation of infinite order, so \ref{rE:I2}, \ref{rE:I3} and \ref{rE:I5} do
not give \ref{rE:I4}.
\end{remark}

\begin{remark}[the exact obstruction]\label{rE:status}
For the Wilson measure with $G=\SU(N)$ and the plaquette sum $\Phi$ the
situation is the following. Everything in Sec.~\ref{ssec:rEunconditional}
holds for every tuning. Of the inputs, \ref{rE:I1} is proved only in the
strong-coupling ranges of Corollaries~\ref{cor:bpranges}
and~\cite[Corollary~322]{BBLR26}, where Theorem~\ref{rE:af}(v) and
Proposition~\ref{rE:necessary}(iv) show that \ref{rE:I2} fails and no
nontrivial limit exists; \ref{rE:I2} is proved in no range of couplings that
tends to the continuum, and by Remark~\ref{rem:limitEafstatus} it is
equivalent to the $Z_k$-free two-sided comparison of
Proposition~\ref{prop:limitEafscale}(iii), open already in a box of fixed
physical size; \ref{rE:I3}, \ref{rE:I4} and \ref{rE:I5} are proved in no
range. By Proposition~\ref{rE:necessary}, (SD) and (NE) with the resulting bound
on the physical gap, which are consequences of \ref{rE:I2} by
Theorem~\ref{rE:af}(iv), are also necessary for (E) with these packets, so
that any proof of (E) along a tuning must establish them; \ref{rE:I2} itself
is not shown to be necessary. The literature supplies none of the five
for the four-dimensional Wilson measure: Ba{\l}aban's renormalisation group
gives bounds on partition functions in a box of fixed physical size uniform
in the spacing~\cite{Bal84a,Bal85a,Bal87,Bal88,Bal89b}, not bounds on the
connected two-point function of $\Phi$; Magnen, Rivasseau and S\'en\'eor
give the basis for a construction of the Schwinger functions of the
$\SU(2)$ gauge field in a regularised axial gauge with an infrared
cut-off~\cite{MRS93}, not expectations of the Wilson measure along a
lattice tuning; the stochastic quantisation of Yang--Mills in two and three
dimensions~\cite{CCHS22,CCHS24} and the state space of the
three-dimensional theory~\cite{CC24} concern lower dimensions; and the weak-coupling correlation decay of
Adhikari and Cao~\cite{AC25} concerns finite gauge groups. The smallest
statement whose proof would change the picture is the two-sided law
\ref{rE:I2} for the $24$-plaquette sum $\Phi$ along the tuning
(AF$_\beta$), whose finite-volume form in a box of fixed physical size is
open as well: for the torus limits along that tuning, Theorem~\ref{rE:af}
would give the renormalisation, the short-distance condition and the
nontriviality of the two-point limits, and with \ref{rE:I1} a two-point
theory with mass gap $m_*$.
\end{remark}
